\newpage
\subsection{Geschäftsleitungsschulung}

\textbf{Anforderung:} Alle Geschäftsleiter müssen regelmäßig zu NIS2, IT-Sicherheit und Risikomanagement geschult werden. Ziel ist, dass sie die Grundlagen von NIS2 verstehen und Risiken für die IT-Sicherheit erkennen und verstehen können.

Zusätzlich zur allgemeinen IT-Sicherheitsschulung müssen Geschäftsleiter regelmäßig speziell zum IT-Risikomanagement geschult werden. Regelmäßig bedeutet auch hier alle drei Jahre. Umstände wie ein Wechsel der Geschäftsleitung oder eine Neuausrichtung der Geschäftsprozesse können eine Abweichung von diesem Zyklus erfordern.

Die Schulung richtet sich primär an Geschäftsleiter, es kann aber sinnvoll sein, auch Personen in ähnlichen Positionen einzubeziehen. Ziel ist der Erwerb von Kenntnissen und Fähigkeiten zur Erkennung und Bewertung von IT-Risiken. Die Schulung kann intern oder extern erfolgen, sofern die Anbieter über die nötigen Kompetenzen verfügen. Wie bei den anderen Schulungen sind Themen, Teilnehmer, Dozent sowie Datum und Dauer als Nachweis zu dokumentieren. Das Format (online oder Präsenz) ist frei, die Dauer kann je nach Vorwissen variieren; das BSI empfiehlt 4 Stunden.

Das BSI (Bundesamt für Sicherheit in der Informationstechnik) hat ein eigenes Dokument zur Geschäftsleitungsschulung veröffentlicht, das detailliert beschreibt, wer geschult werden muss und welche Themen wichtig sind: https://www.bsi.bund.de/dok/nis-2-schulung-geschaeftsleitung
\newpage
Die Schulung für Geschäftsleiter sollte folgende Punkte abdecken:

\begin{itemize}
    \item NIS2-Richtlinie: Grundlegende Informationen über den Zweck, Inhalt und Geltungsbereich der NIS2-Richtlinie \begin{itemize}
        \item Welches Ziel verfolgt die NIS2-Richtlinie?
        \item Wieso ist die eigene Organisation von der NIS2-Richtlinie betroffen? (Branche und Schwellenwerte)
        \item Ist die Organisation nach NIS2 eine wichtige oder besonders wichtige Einrichtung und was sind die Unterschiede?
        \item Welche Themen behandelt die NIS2-Richtlinie und was muss die Organisation umsetzen?
    \end{itemize}
    \item Verantwortung, Haftung und Bußgelder
    \begin{itemize}
        \item Die Geschäftsführer müssen sich bewusst sein, dass sie für die Umsetzung und Einhaltung der NIS2-Richtlinie die Verantwortung tragen und persönlich haften
        \item Welche Bußgelder drohen bei Nichteinhaltung der NIS2-Richtlinie? (für besonders wichtige Einrichtungen bis zu 10 Mio. EUR oder 2\% des weltweiten Jahresumsatzes; für wichtige Einrichtungen bis zu 7 Mio. EUR oder 1,4\% des weltweiten Jahresumsatzes)
        \item Was sind erhebliche Sicherheitsvorfälle und wann müssen diese an das BSI gemeldet werden?
        \item Registrierungspflicht beim BSI von NIS2-betroffenen Organisationen
    \end{itemize}
        \item Umsetzung und Dokumentation
    \begin{itemize}
        \item Welche Rollen, Kompetenzen und Ressourcen werden für die Umsetzung benötigt?
        \item Wie kann ein PDCA-Zyklus implementiert werden und was sind dessen Vorteile?
        \item Was und wie muss für einen Nachweis der Durchführung von NIS2 dokumentiert werden?
    \end{itemize}
    \item Risikomanagement
    \begin{itemize}
        \item Sinn, Ziel und Zweck vom Risikomanagement im Kontext der IT-Sicherheit
        \item Was sind Risiken? (Bedrohung von Geschäftszielen)
        \item Welche Auswirkungen können Vorfälle im Bereich IT-Sicherheit haben? (Malware, Ausfälle bei externen Dienstleistern, Diebstahl von sensiblen Daten, etc.)
        \item Aus welchen Schritten besteht ein Risikomanagement? (Identifizieren, Bewerten, Maßnahmen ergreifen und überwachen)
        \item Welche Frameworks zur Risikoanalyse existieren? (ISO 27005, NIST RMF, etc.)
        \item Was ist die Risikobereitschaft der Organisation? Ab welchen Schwellenwerten müssen bei Risiken Maßnahmen ergriffen werden?
        \item Relevanz von Asset Management für das Risikomanagement
        \item Welche typischen Gefährdungen existieren? (Physische, technische, menschliche Fehler, etc.)
        \item Wie können Risiken identifiziert werden?
        \item Nach welchen Kriterien können Risiken bewertet werden? (Eintrittswahrscheinlichkeit, Schadenshöhe, etc.)
        \item Bewertung von Risiken mittels einer Skala und Abbildung in der Risikomatrix
        \item Welche Maßnahmen können bei Risiken ergriffen werden? (Akzeptieren, Mindern, Vermeiden, Übertragen)
        \item Welche konkreten organisatorischen und technischen Maßnahmen sind möglich?
        \item Wie kann die Wirksamkeit von gewählten Maßnahmen überwacht werden?
        \item Wie kann ein regelmäßiges Risikomanagement in die Prozesse der Organisation eingepflegt werden?
    \end{itemize}
\end{itemize}

\textbf{Checkliste}
\begin{checklist}
\item Prüfen, ob die Schulungen intern oder von externen Anbietern durchgeführt werden
\item Identifizieren, für welche Personen die Schulung verpflichtend ist und für welche sinnvoll
\item Feststellen, über welches Wissen die Geschäftsleiter verfügen und welche Themen in der Schulung behandelt werden müssen
\item Durchgeführte Schulungen (Teilnehmer, Themen, Datum, Dozent) dokumentieren
\item Prüfen und überwachen, unter welchen Bedingungen die Schulung erneut durchgeführt werden muss
\end{checklist}
